% ============================================================
% madXfinder — 추천 알고리즘 (본문: 쉽게 / 부록: 옛 설계에서 바뀐 것)
% Overleaf: New Project → Upload → 이 파일.
% ★ Compiler: XeLaTeX 필수 (Menu → Compiler). pdfLaTeX은 한글 폰트 처리로 시간초과 남
% ============================================================
\documentclass[11pt]{article}
\usepackage{kotex}
\usepackage{amsmath, amssymb}
\usepackage{booktabs}
\usepackage{geometry}
\geometry{a4paper, margin=25mm}
\usepackage{hyperref}
\hypersetup{colorlinks=true, linkcolor=blue, urlcolor=blue}
\usepackage{fancyvrb}
\usepackage{appendix}

\title{madXfinder: 티어표 기반 게임 추천 알고리즘}
\author{madcamp 26s-w1-c3-09}
\date{}

\begin{document}
\maketitle

\section{한 문단 요약}

유저가 즐겨찾기 게임으로 티어표를 만들면, \textbf{``이 게임을 좋아한 사람들이 또 뭘
좋아했나''}(함께 즐겨찾기, co-favorite)를 세서 추천한다. 장르 분류는 안 쓴다 ---
실제 유저들의 즐겨찾기 겹침이 장르보다 취향을 정확히 반영한다는 걸 실험으로 확인했다.
계산에 필요한 데이터는 배치가 미리 DB에 모아두므로, \textbf{일반모드의} 추천 계산은
외부 API를 한 번도 부르지 않는다 (정밀모드는 예외 --- \S\ref{sec:precise}).

\section{점수 계산 (핵심 2줄)}

\subsection{1단계: 티어 가중 합산}

티어표의 게임마다 등급 가중치를 곱해서, 후보 게임의 겹침을 전부 더한다:
\begin{equation}
  \mathrm{raw}(c) = \sum_{g\,\in\,\text{티어표}\cup\text{미배치 즐겨찾기}} w(g)\times\mathrm{overlap}(g,c)
\end{equation}

여기서 $\mathrm{overlap}(g,c)$ = ``$g$를 즐겨찾기한 수집 유저 중 $c$도 즐겨찾기한 사람 수''.

\begin{center}
\begin{tabular}{lcl}
\toprule
등급 & 가중치 $w$ & 비고 \\
\midrule
SSS & 5.5 & 특별 자리 --- 최대 2게임만 \\
A / B / C & 3 / 2 / 1 & \\
미배치 즐겨찾기 & 0.3 & 즐겨찾기했지만 티어에 안 올린 게임 \\
\bottomrule
\end{tabular}
\end{center}

원칙 두 가지 (헷갈리기 쉬운 부분):
\begin{itemize}
  \item \textbf{마이너스 점수는 없다.} 즐겨찾기를 했다는 것 자체가 좋아한다는 뜻이고,
        티어에 안 올린 건 싫어서가 아니라 그냥 안 올린 것이다. 그래서 미배치도 0.3의
        \emph{플러스} 신호.
  \item \textbf{즐겨찾기한 게임은 추천에 안 나온다.} 신호로는 쓰되(입력),
        이미 아는 게임을 추천하면 의미가 없으니 후보에서 전부 뺀다.
\end{itemize}

\subsection{2단계: 유명도 보정}

1단계 점수만 쓰면 초대형 인기 게임이 항상 1등이 된다 (누구의 팬 목록에나 있으니까).
그래서 방문수 $v$로 나눠서 누른다:
\begin{equation}
  \mathrm{score}(c) = \frac{\mathrm{raw}(c)}{v(c)^{\alpha}}
\end{equation}

$\alpha$(알파)가 클수록 유명한 게임이 강하게 눌린다. 하나의 $\alpha$로는
``정확한 추천''과 ``새로운 발견''을 동시에 못 잡아서, \textbf{결과를 두 섹션으로} 낸다:

\begin{center}
\begin{tabular}{lll}
\toprule
섹션 & $\alpha$ & 성격 \\
\midrule
인기 & 0.15 (약하게 누름) & 취향에 맞는 검증된 인기작 \\
발견 & 0.35 (강하게 누름) & 숨어 있는 취향 저격작 \\
\bottomrule
\end{tabular}
\end{center}

마지막으로 필터 둘: \textbf{현재 동접 100명 미만 제외} (문 닫은 게임 거르기 ---
누적 방문수는 과거값이라 못 거르지만 동접은 지금 살아있는지를 정확히 보여줌),
\textbf{겹침 2명 미만 제외} (노이즈 컷).

\subsection{3단계: 나이 보정}

오래된 게임은 팬 데이터가 오래 누적돼 겹침이 과대해지고, 유저 입장에서도 이미
접했을 확률이 높다. 그래서 게임 나이(현재 $-$ 출시일, 년 단위)에 따라 점수에
감쇠 배수를 곱한다 --- 3년까지 1.0, 이후 점차 감소해 9년이면 0.68
(구간 사이는 선형보간, \texttt{scoring.json agePenalty}).

\section{정밀모드: 티어 게임을 그 자리에서 수집}
\label{sec:precise}

일반모드는 배치가 미리 모아둔 cofavorite만 쓰므로, 유저가 티어에 올린 게임이
아직 수집 전이면 그 게임은 추천에 기여하지 못한다. \textbf{정밀모드}는 이를
즉석에서 해결한다 --- 버튼 하나로:
\begin{itemize}
  \item 티어표 \textbf{전체}(SSS/A/B/C, 나이·동접 무관)에 대해 미수집 그룹
        게임의 팬(그룹 멤버 오래된순)을 그 자리에서 수집해 cofavorite를 만들고,
  \item 각 티어 게임의 \textbf{연쇄추천}(로블록스 People-Also-Join, 게임당 6개)을
        ``덤 후보''로 추가한다. 덤 점수 $=$ 최고 소스 티어의 가중치 $\times$ 0.15
        (\texttt{derivedTierFactor}) --- 같은 덤이 여러 티어에서 나오면 최고 티어
        하나만 인정하고, 유저가 직접 배치한 게임이 항상 우선한다.
\end{itemize}
진행률은 게임별 티어 중요도 $\times$ 작업 비용(연쇄추천 1 : 팬수집 표본수)으로
시간에 비례하게 표시하고, 중간에 멈추면 그때까지 수집분으로 즉시 결과를 낸다.
덤 확장을 일반모드에 넣지 않는 이유: 일반모드는 DB에 있는 것만 쓰므로, 연쇄추천
캐시가 있는 게임만 덤을 받아 \textbf{캐시 유무에 따른 쏠림}이 생기기 때문이다.

\section{탐색 범위: depth 1만}

티어표 게임과 \emph{직접} 겹친 게임까지만 본다. ``겹친 게임의 겹침''(depth 2)은
한 다리 건너면서 취향에서 멀어져 노이즈가 되므로 쓰지 않는다 (후보가 부족할 때만 예외).
게임 상세 페이지의 ``함께 즐기는 게임 6개''는 로블록스 공식
People-Also-Join(연쇄추천) --- 캐시에 없으면 그 자리에서 1회 받아와 저장한다.

\section{데이터는 어디서 오나 (요약)}

\begin{Verbatim}[frame=single, fontsize=\small]
배치(Python, 밤):
  인기 게임의 팬(그룹 멤버를 오래된 가입순으로 200명)을 찾아
  → 각자의 즐겨찾기를 전부 DB에 저장
  → "게임 A와 B를 함께 즐겨찾기한 사람 수" 표(co-favorite)를 미리 계산

서버(Java, 실시간):
  유저 요청 → 위 표를 DB에서 읽어 점수 계산 → 응답
  로블록스를 부르는 건 "닉네임→ID 변환"과 "그 유저의 즐겨찾기 1회 조회"뿐
\end{Verbatim}

왜 이렇게 하나: 로블록스 API는 호출 제한이 빡빡해서(예: 게임 상세 초당 1.1회)
실시간으로 부르면 유저 몇 명만 몰려도 밀린다. 자세한 사연은 시행착오 문서 참고.

\section{튜닝은 코드가 아니라 설정 파일에서}

위의 모든 숫자(가중치, $\alpha$, 하한들)는 \texttt{backend/config/scoring.json}에 있다.
값을 바꾸고 서버를 재시작하면 적용된다 --- 코드 수정 불필요.
$\alpha$와 하한들은 실서비스 데이터가 쌓이면 조정할 예정 (미확정 파라미터).

% ============================================================
\appendix
\section*{부록: 처음 설계에서 바뀐 것들}
\addcontentsline{toc}{section}{부록: 처음 설계에서 바뀐 것들}

실측을 거치며 처음 구상이 많이 바뀌었다. 같은 실수를 반복하지 않기 위한 기록.
(상세 근거: \texttt{docs/research/의사결정-기록.md})

\begin{center}
\small
\begin{tabular}{p{0.30\textwidth} p{0.30\textwidth} p{0.30\textwidth}}
\toprule
처음 구상 & 바뀐 것 & 왜 \\
\midrule
로블록스 공식 추천(rec API)을 depth 2까지 연쇄 호출해서 추천
& co-favorite depth 1만
& rec API는 이유를 알 수 없는 블랙박스 + 실시간 호출 예산이 안 나옴.
  co-favorite는 우리가 직접 만든 신호라 설명 가능하고 전부 DB에서 계산됨 \\
\midrule
그룹 멤버 중 일부만 표본 추출 + 뉴비/개발자 스킵 로직
& 오래된 가입순(Asc) 200명, 스킵 없음
& 최신순 뉴비는 즐겨찾기 평균 6.6개로 부실. 오래된순은 95--98\%가 보유(평균 27--30개).
  개발자 스킵은 비교 실험 결과 효과 없음 \\
\midrule
``코어팬'' 판별: 그룹 가입 + 즐겨찾기 수 임계값을 넘는 유저만 저장 (game\_fans 테이블)
& 조회한 유저의 즐겨찾기는 \textbf{무조건 전부 저장}. 팬 = 즐겨찾기한 사람. 테이블도 삭제
& 호출 예산이 귀해서 어렵게 조회한 유저를 버릴 수 없음. 유저 기준으로 쌓으면
  한 번의 조회가 여러 게임 계산에 재사용됨 \\
\midrule
단일 게임 기준 cofav 점수
& 티어표 여러 게임의 \textbf{가중 합산}
& 실서비스 입력은 게임 1개가 아니라 티어표(여러 게임+등급) \\
\midrule
방문수 하한으로 저품질 게임 필터
& \textbf{현재 동접} 하한
& 누적 방문수는 과거값 --- 문 닫은 게임을 못 거름. 동접이 ``지금 살아있는가''를 정확히 반영 \\
\midrule
{[BETA]/[ALPHA]} 이름 필터
& 안 함
& 로블록스는 정식 인기 게임도 그 딱지를 흔히 닮 --- 거르면 진짜 인기작을 놓침 \\
\midrule
장르 다양성 보정
& 안 함
& 같은 유저가 호러·RPG·레이싱을 함께 즐김 --- 장르는 취향의 경계가 아님.
  다양성은 티어표 입력이 알아서 만듦 \\
\bottomrule
\end{tabular}
\end{center}

\end{document}
